\chapter{Conclusion and Future Work}
\label{chap:conclusion}

This project combines customer-focused data warehousing, incremental SSIS integration, versioned K-Means segmentation, and interactive Power BI reporting. The final dimensional design contains two fact tables with complementary purposes: \texttt{FactCustomerRFMLive} provides a purchase-triggered operational customer state, and \texttt{FactCustomerSnapshot} provides comparable month-end observations. The system also retains demographic changes using slowly changing dimensions and stores model parameters so that customer segments can be assigned without repeating training for each purchase.

Historical models were developed for the 2011, 2012, and 2013 training periods, and five business-oriented customer segments were used to interpret the June 2014 dashboard. The dashboards can support practical decisions such as prioritizing a retention offer for \emph{High-Value / Cooling Off} buyers, or designing lower-cost follow-up messages for \emph{Recent / Promising} customers. These recommendations are analytical use cases; the project did not measure the causal effect of a deployed marketing campaign.

The SSIS packages have been deployed for minute-level customer and live RFM refreshes and monthly historical snapshots. The annual Airflow DAG has been configured, but its unattended execution and dependency on the completed December snapshot have not been tested. Further improvement would include completing the Airflow-to-SSIS dependency, reconciling the model and dashboard customer counts, finalizing the Power BI report, validating the source-derived business statistics, and deploying the warehouse online if time and resources allow. These improvements do not change the customer-centred scope of the completed local implementation.
